\chapter{Infraestructura}
\label{cap:infraestructura}

\section{RabbitMQ}

RabbitMQ constituye el bus de integración principal de Atlas y soporta tanto comunicaciones síncronas (RPC) como asíncronas (\textit{Fire-and-Forget}). Su uso transversal permite desacoplar productores y consumidores, reduciendo dependencias temporales entre microservicios.

\subsection*{Topología lógica}

La plataforma opera con tres contextos de mensajería:

\begin{itemize}
\item \textbf{RabbitMQ interno (on-premise)}: tráfico entre microservicios de procesamiento y servicios de IA.
\item \textbf{RabbitMQ IoT + MQTT (on-premise)}: ingestión de eventos de sensores y dispositivos de ciudad.
\item \textbf{RabbitMQ cloud (Azure)}: integración entre API pública, CMS y servicios expuestos en Internet.
\end{itemize}

El \textit{virtual host} de servicio para comunicaciones internas se define como \texttt{/atlas}, aislando permisos y recursos frente a otros contextos del broker.

\subsection*{Patrones de uso}

\begin{itemize}
\item \textbf{RPC}: una petición publica en una cola de servicio y espera respuesta mediante \texttt{reply\_to} y \texttt{correlation\_id}.
\item \textbf{Fire-and-Forget}: publicación de eventos de dominio sin bloqueo del emisor.
\item \textbf{Enrutado por etiquetas}: los \textit{tags} canónicos de la ontología (\texttt{dominio:categoría:acción[:subtipo]}) se transforman a \textit{routing keys} AMQP mediante sustitución \texttt{:} $\rightarrow$ \texttt{.}.
\end{itemize}

\subsection*{Prácticas de configuración}

\begin{itemize}
\item Exchanges y colas declarados por servicio durante inicialización.
\item Usuarios técnicos separados por rol (emisor, receptor, administración).
\item Políticas de durabilidad y reintento adaptadas al tipo de mensaje.
\item Restricción de permisos por vhost, exchange y patrón de routing key.
\end{itemize}

\section{Base de datos (MariaDB)}

MariaDB proporciona persistencia relacional para entidades de negocio, configuración editorial y metadatos de integración. En Atlas se despliega en ambos entornos (on-premise y cloud) para mantener autonomía operativa.

\subsection*{Modelo de uso}

\begin{itemize}
\item Almacenamiento transaccional de contenidos estructurados del CMS.
\item Persistencia de relaciones entre entidades de dominio y taxonomías.
\item Soporte de consultas indexadas para consumo API y procesos batch.
\end{itemize}

\subsection*{Conectividad y operación}

\begin{itemize}
\item Puerto de servicio: \texttt{3306}.
\item Acceso mediante credenciales por aplicación y principio de mínimo privilegio.
\item Gestión de conexiones desde librerías de infraestructura (\textit{pooling} y reutilización).
\item Copias de seguridad y políticas de retención definidas por entorno.
\end{itemize}

Las operaciones de escritura intensiva se desacoplan del plano de exposición pública mediante colas de eventos, reduciendo picos de carga directa sobre la base de datos.

\section{Caché (Valkey)}

Valkey se utiliza como capa de caché en memoria compatible con Redis para reducir latencias de lectura y descargar carga de MariaDB en operaciones repetitivas.

\subsection*{Casos de uso}

\begin{itemize}
\item Caché de consultas frecuentes del CMS y de respuestas de lectura.
\item Almacenamiento temporal de sesiones, tokens y artefactos de corta vida.
\item Mecanismos de invalidación coordinada tras cambios editoriales.
\end{itemize}

\subsection*{Parámetros operativos}

\begin{itemize}
\item Puerto de servicio: \texttt{6379}.
\item Estrategia de expiración basada en TTL por tipo de dato.
\item Segmentación por prefijos de clave para evitar colisiones entre servicios.
\item Políticas de memoria ajustadas al perfil de cada entorno.
\end{itemize}

\section{Modelos de IA (Ollama)}

El procesamiento de inteligencia artificial se ejecuta principalmente en el CPD municipal mediante Ollama, aprovechando capacidad GPU local y manteniendo soberanía del dato para cargas sensibles.

\subsection*{Capacidades funcionales}

\begin{itemize}
\item \textbf{LLM}: clasificación, resumen y extracción semántica de contenidos.
\item \textbf{OCR}: transcripción de documentos e imágenes para su normalización posterior.
\item \textbf{Embeddings}: vectorización de textos para búsqueda semántica y similitud.
\end{itemize}

\subsection*{Modelos y conectividad}

\begin{itemize}
\item Endpoint de servicio: \texttt{11434}.
\item Modelos base operativos: \texttt{gemma3}, \texttt{deepseek-ocr}, modelos de embeddings compatibles.
\item Integración vía microservicio \texttt{atlas-service-ollama} y mensajería AMQP.
\end{itemize}

La invocación de capacidades IA se encapsula en contratos de mensajería definidos en \textit{atlas-messaging-ontology}, permitiendo sustitución de modelos sin romper interfaces funcionales.

\section{Proxy inverso (NGINX)}

NGINX actúa como capa de entrada HTTP/HTTPS en ambos entornos y centraliza funciones de publicación, seguridad y encaminamiento hacia los servicios internos.

\subsection*{Funciones principales}

\begin{itemize}
\item Terminación TLS y redirección obligatoria de HTTP a HTTPS.
\item Enrutado por host/path hacia API Gateway, CMS y servicios auxiliares.
\item Aplicación de cabeceras de seguridad y límites de tamaño de petición.
\item Registro de accesos para auditoría y diagnóstico operativo.
\end{itemize}

\subsection*{Puertos y exposición}

\begin{itemize}
\item \texttt{80/tcp}: tráfico HTTP y redirección controlada.
\item \texttt{443/tcp}: tráfico HTTPS productivo.
\end{itemize}

La configuración del proxy se versiona por entorno para separar certificados, dominios públicos y reglas de encaminamiento internas.
